\chapter{Metodologi Penelitian}

\section{Alur Penelitian}
\TemplateTodo{Jelaskan tahapan penelitian dari awal hingga akhir. Pastikan alur sesuai dengan rumusan masalah dan hasil yang akan dilaporkan.}

Metodologi harus cukup rinci agar penelitian dapat dipahami dan, sejauh mungkin, direplikasi. Tuliskan keputusan penting, parameter, asumsi, serta alasan pemilihan metode.

\section{Objek dan Data Penelitian}
\TemplateTodo{Jelaskan objek penelitian, sumber data, kriteria inklusi/eksklusi, periode pengambilan data, jumlah data, dan aspek etik atau perizinan jika relevan.}

\section{Analisis Kebutuhan}
\TemplateTodo{Uraikan kebutuhan perangkat keras, perangkat lunak, dan kebutuhan fungsional sistem. Sebutkan mengapa setiap kebutuhan membatasi pilihan yang diambil pada bab ini.}

Spesifikasi lingkungan yang digunakan ditunjukkan pada Tabel \ref{tab:kebutuhan-perangkat}.

% CONTOH TABEL 2: sebagian kolom lebarnya tetap, satu kolom terakhir lentur.
%
% Bentuk inilah yang paling sering dipakai. Kolom P{0,20\linewidth} lebarnya
% dikunci, sedangkan kolom Y menyerap sisa lebar halaman. Letakkan kolom yang
% isinya paling panjang di posisi Y, sehingga kolom pendek tidak diregangkan
% sampai berlubang.
%
% Jumlah lebar seluruh kolom P sebaiknya di bawah 1,0\linewidth agar kolom Y
% masih kebagian ruang. Di bawah ini 0,20 + 0,30 = 0,50, menyisakan sekitar
% setengah lebar halaman untuk kolom keterangan.
\begin{xltabular}{\linewidth}{|P{0.20\linewidth}|P{0.30\linewidth}|Y|}
\caption{Spesifikasi lingkungan komputasi yang digunakan.}\label{tab:kebutuhan-perangkat}\\
\hline
Komponen & Spesifikasi & Keterangan \\
\hline
\endfirsthead
\hline
Komponen & Spesifikasi & Keterangan \\
\hline
\endhead
\hline
\endlastfoot
Prosesor & Sebutkan model dan jumlah inti & Jelaskan pengaruhnya terhadap konfigurasi yang dipilih \\
\hline
Memori & Sebutkan kapasitas RAM & Jelaskan batas ukuran data yang dapat diproses sekaligus \\
\hline
Perangkat lunak & Sebutkan bahasa, pustaka, dan versinya & Versi dicantumkan agar hasil dapat direplikasi \\
\end{xltabular}

Dua batasan pada Tabel \ref{tab:kebutuhan-perangkat} berdampak langsung pada konfigurasi penelitian. Jelaskan dampaknya di sini, jangan sekadar mengulang isi tabel.

\section{Perancangan Metode atau Sistem}
\TemplateTodo{Uraikan rancangan metode, model, sistem, atau prosedur eksperimen. Hubungkan setiap komponen dengan tujuan penelitian.}

% \Placeholder menandai bagian yang belum ditulis. Kotaknya tampil pada
% main.tex, tetapi hilang pada proposal.tex dan skripsi-draft.tex, sehingga
% naskah yang dibawa ke bimbingan tetap bersih. Atur \ThesisShowHelpers di
% config/thesis-config.tex bila ingin memaksa tampil atau selalu disembunyikan.
\Placeholder{Rincian arsitektur sistem masih menunggu hasil uji coba awal.}

% \FigurePlaceholder mengisi tempat gambar yang belum tersedia. Berbeda dengan
% mengosongkan seluruh environment figure, cara ini membuat nomor gambar,
% caption, dan \label tetap hidup, jadi rujukan \ref di paragraf lain tidak
% berubah ketika gambar aslinya nanti dipasang.
\begin{figure}[H]
\centering
\FigurePlaceholder{4cm}{Diagram rancangan sistem yang diusulkan}
\caption{Rancangan sistem yang diusulkan}
\label{fig:rancangan-sistem}
\end{figure}

Rancangan sistem ditunjukkan pada Gambar~\ref{fig:rancangan-sistem}. Ganti \verb|\FigurePlaceholder| dengan \verb|\includegraphics| begitu gambarnya siap.

\section{Skenario Pengujian}
\TemplateTodo{Jelaskan rancangan eksperimen, pembagian data, pembanding, metrik evaluasi, perangkat yang digunakan, dan cara memastikan hasil adil.}

\section{Teknik Analisis Data}
\TemplateTodo{Jelaskan bagaimana hasil akan dianalisis: statistik deskriptif, perbandingan metrik, analisis kesalahan, validasi pengguna, atau pendekatan lain yang sesuai.}
